\documentclass[a4paper,12pt,oneside]{report}
\usepackage{graphicx}
\usepackage{color}
\usepackage{url}
\usepackage{subfigure}
\usepackage[utf8]{inputenc}
\usepackage[T1]{fontenc}
%\usepackage{tgpagella}
\usepackage{xstring}
\usepackage{wrapfig}
\usepackage{appendix}

\newcommand{\myscale}{0.74}
\newcommand{\vect}[1]{\boldsymbol{#1}}
\newcommand{\code}[1]{\texttt{\StrSubstitute{#1}{.}{.\.}}}
\def\.{\discretionary{}{}{}}
\newcommand{\jmodule}[1]{\texttt{\textit{#1}}}
\linespread{1.5}

\setlength{\hoffset}{-1in} %left margin will be 0, as hoffset is by default 1inch
\setlength{\voffset}{-1in} %analogous voffset
\setlength{\oddsidemargin}{4cm}
\setlength{\evensidemargin}{4cm}
\setlength{\topmargin}{25mm}
\setlength{\footskip}{1cm}
\setlength{\headheight}{0cm}
\setlength{\headsep}{0cm}
\setlength{\marginparwidth}{0cm}
\setlength{\marginparpush}{0cm}
\setlength{\textheight}{23.7cm}
\setlength{\textwidth}{14.5cm}
\let\openright=\clearpage

\def\mfauthor{Matej Vitásek}
\def\mfadvisor{RNDr. Irena Mlýnková Ph.D.}
\def\mfplacedate{Praha, 2011}

% Tato makra přesvědčují mírně ošklivým trikem LaTeX, aby hlavičky kapitol
% sázel příčetněji a nevynechával nad nimi spoustu místa. Směle ignorujte.
\makeatletter
\def\@makechapterhead#1{
  {\parindent \z@ \raggedright \normalfont
   \Huge\bfseries \thechapter. #1
   \par\nobreak
   \vskip 20\p@
}}
\def\@makeschapterhead#1{
  {\parindent \z@ \raggedright \normalfont
   \Huge\bfseries #1
   \par\nobreak
   \vskip 20\p@
}}
\makeatother

\def\chapwithtoc#1{
\chapter*{#1}
\addcontentsline{toc}{chapter}{#1}
}

\begin{document}

\lefthyphenmin=2
\righthyphenmin=2

%%% Titulní strana práce ======================================================
\pagestyle{empty}
\begin{center}

\large

Charles University in Prague

\medskip

Faculty of Mathematics and Physics

\vfill

{\bf\Large MASTER THESIS}

\vfill

\centerline{\mbox{\includegraphics[width=60mm]{logo.eps}}}

\vfill
\vspace{5mm}

{\LARGE \mfauthor}

\vspace{15mm}

% exactly as assigned
{\LARGE\bfseries Inference of XML Integrity Constraints}

\vfill

% Název katedry nebo ústavu, kde byla práce oficiálně zadána
% (dle Organizační struktury MFF UK)
Department of Software Engineering

\vfill

\begin{tabular}{rl}

Supervisor of the master thesis: & 	\mfadvisor \\
\noalign{\vspace{2mm}}
Study programme: & Informatika \\
\noalign{\vspace{2mm}}
Specialization: & ISS \\
\end{tabular}

\vfill

% Zde doplňte rok
\mfplacedate 

\end{center}



\newpage % ============================================================
%%% Následuje vevázaný list -- kopie podepsaného "Zadání diplomové práce".
%%% Toto zadání NENÍ součástí elektronické verze práce, nescanovat.
\openright

\noindent
[Sample: Here you may thank whoever you wish (the supervisor of the thesis, the consultant,
the person who lent the software, literature, etc.)]

TODO thank jInfer team, Mlynkova, anyone helping create this, colleagues at work...


\newpage % ============================================================
%%% Strana s čestným prohlášením k diplomové práci
\vglue 0pt plus 1fill
\noindent
I declare that I carried out this master thesis independently, and only with the cited
sources, literature and other professional sources.

\medskip\noindent
I understand that my work relates to the rights and obligations under the Act No.
121/2000 Coll., the Copyright Act, as amended, in particular the fact that the Charles
University in Prague has the right to conclude a license agreement on the use of this
work as a school work pursuant to Section 60 paragraph 1 of the Copyright Act.

\vspace{10mm}

\hbox{\hbox to 0.5\hsize{%
In ........ date ............
\hss}\hbox to 0.5\hsize{%
signature
\hss}}

\vspace{20mm}



\newpage % ============================================================
%%% Povinná informační strana diplomové práce
\vbox to 0.5\vsize{
\setlength\parindent{0mm}
\setlength\parskip{5mm}

Název práce: Title [accurate translation of the master thesis title in Czech]

Autor: \mfauthor

Katedra:  Katedra softwarového inženýrství

Vedoucí diplomové práce: \mfadvisor{}, Katedra softwarového inženýrství

Abstrakt: [abstract of 80-200 words in Czech, but not a copy of the assignment of the
master thesis]

Klíčová slova:  [3-5 keywords in Czech]

\vss}\nobreak\vbox to 0.49\vsize{
\setlength\parindent{0mm}
\setlength\parskip{5mm}

Title:  Inference of XML Integrity Constraints

Author: \mfauthor

Department: Department of Software Engineering

Supervisor: \mfadvisor{}, Department of Software Engineering

Abstract:  [abstract of 80-200 words in English, but not a copy of the assignment of
the master thesis]

Keywords: XML, ID attributes, inference

\vss}

\newpage

%%% Strana s automaticky generovaným obsahem diplomové práce. U matematických
%%% prací je přípustné, aby seznam tabulek a zkratek, existují-li, byl umístěn
%%% na začátku práce, místo na jejím konci.

\openright
\pagestyle{plain}
\setcounter{page}{1}
\tableofcontents

% =============== TEXT ================
\newpage

\chapwithtoc{Preface}

TODO why is this important

\chapter{Definitions}

TODO talk about key definition from \cite{keX}

talk about ID, IDREF, ... talk about XML schema keys

\chapter{Research}

According to the article \cite{fidax}, ...

TODO talk about FIDAX
To the best of our knowledge, there are no other articles dealing with this problem 

TODO talk about Fajt - but that's different, that's keys

\section{Independent Set}

TODO other approaches to max IS

\chapter{MIP Approach}

TODO FIDAX is cool, but their heuristic might not be the best. It is a IS problem, let's try a multi-purpose solver.

\section{Finding ID sets with GLPK}

TODO how GLPK works, how to construct a problem, how output looks, how to interpret it, limit time-> limit quality, ...
 
 -> wow, GLPK works. However, it takes too long to get to optimum (link experiments), so let's try other heuristics.
 
\section{Heuristics}

TODO what is heuristics (link wise books), what is metaheuristics (we will be using them)
TODO mention things like Taboo search

TODO we will be working with heus striving to do the following: input is a list of AMs, they have their weight, we try to find a non-conflicting subset which maximizes the weight

\subsection{Constructions Heuristics}

TODO construction heuristics are heus that provide us with at least some solution.

\subsubsection{FIDAX}

TODO this is a heuristic too! We can start with it!

\subsubsection{Random}

TODO this is trivial, but works - link experiments where it bashes FIDAX

\subsubsection{Fuzzy}

TODO kind of trivial

\subsubsection{Incremental}

TODO trivial, hungry

\subsubsection{Removal}

TODO trivial, "hungry"

\subsubsection{Truncated Branch \& Bound}

TODO if we limit GLPK's runtime, we get this 

\subsection{Improvement Heuristics}

TODO what they are, that they need a pool sometimes, how do we give them pool, ...

\subsubsection{Identity}

TODO trivial trivial

\subsubsection{Remove Worst}

TODO trivial

\subsubsection{Random Remove}

TODO trivial, perturbation

\subsubsection{Hungry}

TODO trivial

\subsubsection{Mutation}

TODO

\subsubsection{Crossover}

TODO

\subsubsection{Local Branching}

TODO

\chapter{Experiments}

\section{Experimental Data}

\section{Experimental Setup}

\section{Experimental Results}

TODO following experiments + more

Input -> How long the search for optimum lasts
Time limit -> Quality
Structure of the data -> Which algorithms perform the best
Alpha, Beta -> Quality
Which GLPK -> *
Do I beat FIDAX with random?
Does Fuzzy beat Random? (Seems that Random is better and even has smaller variance)

\section{The "Best" Algorithm}

\chapwithtoc{Conclusion}

TODO conclusion

%%% Seznam použité literatury
\newpage
\nocite{*}
\bibliographystyle{alpha}
\bibliography{literature}

%%% Tabulky v diplomové práci, existují-li.
%\chapwithtoc{List of Tables}
\listoftables

%%% Použité zkratky v diplomové práci, existují-li, včetně jejich vysvětlení.
\chapwithtoc{List of Abbreviations}

%%% Přílohy k diplomové práci, existují-li (různé dodatky jako výpisy programů,
%%% diagramy apod.). Každá příloha musí být alespoň jednou odkazována z vlastního
%%% textu práce. Přílohy se číslují.

%% TODO Appendix - About jInfer 

\appendix

\openright
\addcontentsline{toc}{chapter}{Appendices}

\chapter{jInfer}

\end{document}
